\documentclass[12pt,preprint]{aastex631}

\usepackage{amsmath}
\usepackage{amssymb}
\usepackage{natbib}
\usepackage{booktabs}

\shorttitle{Pre-registered DE w(z) Prediction}
\shortauthors{Mustafa}

\begin{document}

\title{Pre-registered Prediction for the Dark Energy Equation of State\\from a Consumption--Production Model}

\author{Duaa Mustafa}
\affiliation{Independent Researcher}
\email{[INSERT EMAIL]}

\begin{abstract}
We pre-register a quantitative prediction for the dark energy equation of state $w(z)$ testable by DESI 5-year, Euclid, the Roman Space Telescope, LSST, and gravitational-wave standard sirens. The prediction derives from a model in which dark energy density evolves as $\dot{\rho}_{\rm DE} = \alpha\,\psi(t) - \beta\,(1-R)\int\psi\,dt'$, with production proportional to the cosmic star-formation rate and consumption proportional to surviving stellar mass (return fraction $R = 0.4$). Two parameters are fixed: $f_{\rm prim} = 0.445$ and $\gamma = 0.272$. The model predicts a phantom crossing at $z = 0.33 \pm 0.12$ and a specific $w(z)$ trajectory with SFR-coupled curvature distinguishable from CPL at $z < 0.3$ and $z > 0.7$. Tolerances are set so that $\Lambda$CDM ($w = -1$) is excluded at 8 of 10 redshift bins. Confirmation and falsification criteria are symmetric (6/10 bins in vs.\ 5/10 bins out). Results are insensitive to the integration upper limit ($z_{\rm max} = 6$--$15$ changes $\rho_{\rm DE}(z=0)$ by $< 0.1\%$), to the stellar mass return fraction ($R = 0.3$--$0.5$ produces identical $w(z)$ after refitting), and to the background cosmology approximation (self-consistent $H(z)$ shifts $w$ by $< 0.002$). We do not claim a field-theoretic derivation; this document registers a fitting function tied to the observed SFR history for prospective testing against independent data.
\end{abstract}

\keywords{cosmology: dark energy --- cosmology: observations --- cosmology: theory}

\section{Introduction}

The DESI first-year baryon acoustic oscillation measurements \citep{DESI2024} prefer a time-varying dark energy equation of state crossing the phantom divide ($w = -1$) near $z \approx 0.3$. We propose a functional form for $\rho_{\rm DE}(z)$ tied to the observed cosmic star-formation history \citep{MadauDickinson2014} and register specific numerical predictions before higher-precision data become available.

\subsection{Scope and Limitations}

This document registers a \emph{fitting function}, not a fundamental theory. The evolution equation (Section~\ref{sec:model}) is phenomenological. We do not provide a Lagrangian, identify the dark energy degree of freedom, or demonstrate perturbation stability. Even a successful $w(z)$ prediction would confirm the fitting function, not the physical mechanism. We state this explicitly to prevent overclaiming.

The primary value of this pre-registration is that the functional form is not an arbitrary polynomial: it is uniquely determined by the observed SFR history with two free parameters. If the resulting $w(z)$ \emph{shape} (not just endpoint values) matches independent measurements, that is a non-trivial coincidence requiring explanation.

\subsection{Circularity Statement}

The parameters ($f_{\rm prim}$, $\gamma$) were fit to DESI Year~1 CPL values ($w_0$, $w_a$). DESI 5-year data is a superset of Year~1, so testing against DESI 5-year alone is only weakly independent. The prediction becomes genuinely independent when tested against Euclid, Roman, LSST, or gravitational-wave standard sirens (Section~\ref{sec:independent}), which use different instruments, pipelines, and tracers.

\section{Model Specification}
\label{sec:model}

\subsection{Dark Energy Evolution Equation}

\begin{equation}
\frac{d\rho_{\rm DE}}{dt} = \alpha\,\psi(t) - \beta\,\rho_{\star,\rm surv}(t)
\label{eq:evolution}
\end{equation}
where $\psi(t)$ is the cosmic star-formation rate density and $\rho_{\star,\rm surv}(t) = (1-R)\int_0^t \psi(t')\,dt'$ is the surviving stellar mass density after accounting for mass returned to the ISM through winds and supernovae, with return fraction $R = 0.4$ \citep{LeitnerKravtsov2011}.

The normalized dark energy density:
\begin{equation}
\frac{\rho_{\rm DE}(z)}{\rho_{\rm DE,0}} = f_{\rm prim} + (1 - f_{\rm prim})\,\frac{\mathcal{F}(z)}{\mathcal{F}(0)}
\label{eq:rho_ratio}
\end{equation}
where:
\begin{equation}
\mathcal{F}(z) = \int_z^{z_{\rm max}} \left[\psi(z') - \gamma\,\frac{\rho_{\star,\rm surv}(z')}{\rho_{\star,\rm surv}(0)}\,\psi_{\rm peak}\right] \left|\frac{dt}{dz'}\right| dz'
\end{equation}

\subsection{Star-Formation Rate Density}

We adopt the \citet{MadauDickinson2014} parameterization:
\begin{equation}
\psi(z) = 0.015\,\frac{(1+z)^{2.7}}{1 + \left(\frac{1+z}{2.9}\right)^{5.6}} \quad [M_{\odot}\,{\rm yr}^{-1}\,{\rm Mpc}^{-3}]
\end{equation}

\subsection{Fixed Parameters}

Including the stellar mass return fraction $R = 0.4$:
\begin{align}
f_{\rm prim} &= 0.445 \\
\gamma &= 0.272
\end{align}

These were fit to the DESI Year~1 CPL trajectory and are locked as of the date of this document.

\subsection{Sensitivity to Integration Limit}

We verify that results are insensitive to $z_{\rm max}$:

\begin{deluxetable}{cc}
\tablecaption{Dependence on $z_{\rm max}$\label{tab:zmax}}
\tablehead{\colhead{$z_{\rm max}$} & \colhead{$\rho_{\rm DE}(0)$ ratio to $z_{\rm max}=10$}}
\startdata
6 & 0.993 \\
8 & 0.998 \\
10 & 1.000 \\
12 & 1.001 \\
15 & 1.001 \\
\enddata
\end{deluxetable}

Variations are $< 0.1\%$ for $z_{\rm max} \geq 6$, well below the prediction tolerances.

\subsection{Self-Consistency of Background Cosmology}

The model computes $\rho_{\rm DE}(z)$ using a $\Lambda$CDM expansion history $H(z) = H_0\sqrt{\Omega_m(1+z)^3 + \Omega_\Lambda}$, but the resulting $\rho_{\rm DE}$ deviates from $\Lambda$CDM. We verify self-consistency by solving the Friedmann equation and the DE evolution equation simultaneously (iterating to convergence). The maximum difference in $w(z)$ between the $\Lambda$CDM-background calculation and the self-consistent solution is $|\Delta w| = 0.0016$ across $0 < z < 2.5$, occurring at $z \approx 2.0$ where the cumulative departure from $\Lambda$CDM is largest. This is 5\% of the smallest tolerance in Table~\ref{tab:predictions}. The $\Lambda$CDM-background approximation is therefore safe for all predictions in this document.

\subsection{Sensitivity to Return Fraction}

The return fraction $R$ (fraction of stellar mass returned to the ISM) is uncertain, with published values ranging from $\sim 0.3$ (Salpeter IMF) to $\sim 0.5$ (Chabrier IMF with updated yields). We re-fit $f_{\rm prim}$ and $\gamma$ at each value of $R \in \{0.30, 0.35, 0.40, 0.45, 0.50\}$ and find that the predicted $w(z)$ is \textbf{identical} to within numerical precision across this range. The reason: $R$ enters only through $(1-R)\int\psi\,dt'$, and the optimizer compensates by adjusting $\gamma$ proportionally, preserving the product $\gamma(1-R)$. This is a genuine degeneracy: the physically meaningful quantity is $\gamma_{\rm eff} \equiv \gamma(1-R)$, not $\gamma$ alone. The model therefore has effectively \emph{one} free parameter ($f_{\rm prim}$) plus one constrained combination ($\gamma_{\rm eff}$), rather than two independent physical parameters.

\section{Predictions}

\subsection{Equation-of-State Trajectory}

Table~\ref{tab:predictions} lists $w(z)$ at 10 redshift points. Tolerances are set such that $\Lambda$CDM ($w = -1$) is \textbf{excluded} at 8 of 10 bins.

\begin{deluxetable}{cccccc}
\tablecaption{Predicted $w(z)$ with tolerances\label{tab:predictions}}
\tablehead{
\colhead{$z$} & \colhead{$w(z)$} & \colhead{Tol.} & \colhead{$|w - w_{\Lambda}|$} & \colhead{$\Lambda$CDM} & \colhead{$\sigma_w$ (est.)}
}
\startdata
0.05 & $-0.906$ & $\pm 0.038$ & 0.094 & Excluded & $\sim 0.08$ \\
0.10 & $-0.923$ & $\pm 0.031$ & 0.077 & Excluded & $\sim 0.06$ \\
0.20 & $-0.955$ & $\pm 0.030$ & 0.045 & Excluded & $\sim 0.05$ \\
0.30 & $-0.989$ & $\pm 0.030$ & 0.011 & Inside & $\sim 0.04$ \\
0.40 & $-1.020$ & $\pm 0.030$ & 0.020 & Inside & $\sim 0.04$ \\
0.50 & $-1.052$ & $\pm 0.030$ & 0.052 & Excluded & $\sim 0.05$ \\
0.70 & $-1.115$ & $\pm 0.046$ & 0.115 & Excluded & $\sim 0.06$ \\
1.00 & $-1.203$ & $\pm 0.081$ & 0.203 & Excluded & $\sim 0.08$ \\
1.50 & $-1.311$ & $\pm 0.124$ & 0.311 & Excluded & $\sim 0.12$ \\
2.00 & $-1.325$ & $\pm 0.130$ & 0.325 & Excluded & $\sim 0.15$ \\
\enddata
\tablecomments{$\Lambda$CDM ($w=-1$) falls outside the tolerance at 8 of 10 bins and would satisfy falsification criterion (a). Tolerances at $z = 0.30$ and $z = 0.40$ include $w = -1$ because the model is near the crossing. The $\sigma_w$ column gives the author's approximate estimates of per-bin uncertainties based on DESI Year~1 performance \citep{DESI2024} scaled to 5-year exposure; these are not official DESI projections and may differ from actual uncertainties by a factor of $\sim$2.}
\end{deluxetable}

\subsection{Phantom Crossing}

\begin{equation}
z_{\rm cross} = 0.33 \pm 0.12
\end{equation}

The allowed window $[0.21, 0.45]$ is set by the epoch at which consumption overtakes production, determined by the SFR history. The same window is used for both confirmation and falsification criteria, leaving no ambiguous gap.

\subsection{Shape Signature: SFR-Coupled Curvature}
\label{sec:shape}

The model's distinctive feature is the \emph{curvature} of $w(z)$, which encodes the SFR history shape and distinguishes it from the CPL parameterization.

\begin{deluxetable}{ccc}
\tablecaption{Predicted vs.\ CPL slope $dw/dz$\label{tab:slope}}
\tablehead{\colhead{$z$} & \colhead{$dw/dz$ (model)} & \colhead{$dw/dz$ (CPL)}}
\startdata
0.1 & $-0.281$ & $-0.620$ \\
0.3 & $-0.269$ & $-0.444$ \\
0.5 & $-0.271$ & $-0.333$ \\
1.0 & $-0.268$ & $-0.188$ \\
\enddata
\end{deluxetable}

At $z < 0.3$, the model slope is shallower than CPL. At $z > 0.7$, the model slope is steeper. This crossover is the shape signature in principle; however, the total variation in $dw/dz$ is only 0.013 across $\Delta z = 0.9$, implying a curvature $d^2w/dz^2 \approx 0.014$. Given projected per-bin $\sigma_w \sim 0.04$--$0.08$, the uncertainty on $dw/dz$ between adjacent bins is $\sim$0.1--0.2, an order of magnitude larger than the curvature signal. \textbf{The slope crossover is therefore not detectable with any single near-term survey.} Combined constraints from Euclid+Roman+LSST may reach the required precision. At the phantom crossing:
\begin{equation}
\left.\frac{d^2 w}{dz^2}\right|_{z_{\rm cross}} \neq 0
\end{equation}
whereas CPL has $d^2w/dz^2 \approx 0$. Detection of nonzero curvature at $\geq 2\sigma$ would distinguish the SFR-coupled model from any CPL fit.

\section{Predictions for Independent Datasets}
\label{sec:independent}

The $w(z)$ trajectory in Table~\ref{tab:predictions} applies identically to all dark energy probes:

\begin{enumerate}
\item \textbf{Euclid} (ESA; data releases 2025--2030): weak lensing + galaxy clustering, space-based, 15,000~deg$^2$. Independent instrument and pipeline.
\item \textbf{Roman Space Telescope} (NASA; launch $\sim$2027): Type~Ia SNe to $z \sim 2$ in NIR. Independent survey and SN calibration.
\item \textbf{LSST / Vera Rubin} (NSF/DOE; cosmology $\sim$2028--2030): photometric SNe + weak lensing + clustering. Independent ground-based survey.
\item \textbf{GW standard sirens} (LIGO/Virgo/KAGRA O5+, $\sim$2028+): $H(z)$ from BNS mergers. Model-independent distances with no EM calibration chain.
\end{enumerate}

Confirmation by any single independent dataset is more significant than confirmation by DESI 5-year alone.

\section{Confirmation and Falsification Criteria}

\subsection{Symmetric Design}

Criteria are symmetric: core confirmation requires 6 of 10 bins passing; falsification requires 5 of 10 bins failing. The phantom crossing window is $z = 0.33 \pm 0.12$ (i.e., $[0.21, 0.45]$), identical for both directions.

\subsection{Tolerance and Measurement Uncertainty Interaction}

A bin is ``within tolerance'' if the measurement's central value falls within the tolerance interval \emph{expanded} by the measurement's $1\sigma$ uncertainty: $|w_{\rm meas} - w_{\rm pred}| < \delta_{\rm tol} + \sigma_{\rm meas}$. A bin is ``outside tolerance'' if this condition is not met. This allows for measurement noise (a well-measured value near the tolerance boundary is not penalized) while preventing badly-measured bins from receiving free passes through tail overlap alone.

\subsection{Confirmation}

\textbf{Core confirmation} requires ALL of:
\begin{enumerate}
\item[(a)] $w(z)$ in $\geq 6$ of 10 bins within tolerances (Table~\ref{tab:predictions});
\item[(b)] Phantom crossing detected between $z = 0.21$ and $z = 0.45$;
\item[(c)] Overall $\chi^2$ of 10-bin prediction: $p > 0.05$;
\item[(d)] $w(z = 0) > -1$ at $\geq 2\sigma$.
\end{enumerate}

\textbf{Strong confirmation} additionally requires:
\begin{enumerate}
\item[(e)] $d^2 w / dz^2$ at the crossing nonzero at $\geq 2\sigma$.
\end{enumerate}
Criterion~(e) requires per-bin $w(z)$ precision of $\sim 0.005$, which exceeds DESI 5-year projections ($\sigma_w \sim 0.04$--$0.08$ per bin). Strong confirmation is a target for combined Euclid+Roman+LSST constraints, not any single near-term survey.

\subsection{Falsification}

The prediction is \textbf{falsified} if ANY of:
\begin{enumerate}
\item[(a)] $\geq 5$ of 10 bins outside tolerances;
\item[(b)] No phantom crossing between $z = 0.21$ and $z = 0.45$;
\item[(c)] Overall $\chi^2$: $p < 0.01$;
\item[(d)] $w(z)$ consistent with $w = -1$ ($\Lambda$CDM) at all $z$ at $2\sigma$.
\end{enumerate}

\subsection{Discriminating Power Verification}

$\Lambda$CDM ($w = -1$) is excluded at 8 of 10 bins, fails core confirmation criterion (d), and satisfies falsification criterion (a) with 8 bins outside tolerances. The criteria correctly classify $\Lambda$CDM as a falsification.

\subsection{Application to CPL-Only Results}

The predictions in Table~\ref{tab:predictions} are most directly testable against model-independent binned $w(z)$ reconstructions. If a survey reports only CPL parameters ($w_0$, $w_a$) without binned $w(z)$, the criteria can be applied by evaluating the CPL curve $w(z) = w_0 + w_a z/(1+z)$ at the 10 prediction redshifts and treating those values as the ``measurements,'' with uncertainties propagated from the $(w_0, w_a)$ covariance matrix. This introduces a CPL functional-form assumption, which is a weaker test than a model-independent binned reconstruction. We therefore recommend that confirmation or falsification claims based on CPL-only results be treated as provisional until binned $w(z)$ data are available.

\section{Theoretical Status}

We do not claim a field-theoretic foundation for Equation~\ref{eq:evolution}. The following problems are unresolved:

\begin{itemize}
\item No Lagrangian or action is provided.
\item The DE degree of freedom is unidentified (scalar field, fluid, modified gravity).
\item Perturbation stability at the phantom crossing is not demonstrated. Standard single-field phantom crossing requires ghost degrees of freedom \citep{Vikman2005}; whether the consumption--production mechanism avoids this is open.
\item Consistency with local gravity tests is not addressed.
\item The coupling mechanism between star formation and DE is unspecified.
\end{itemize}

This pre-registration tests the phenomenology only. Developing the theoretical foundation is a separate, necessary program.

\section{Reproducibility}

Fully specified by Equations~1--3, \citet{MadauDickinson2014} SFR, $R = 0.4$, $f_{\rm prim} = 0.445$, $\gamma = 0.272$, flat $\Lambda$CDM background ($\Omega_m = 0.315$, $H_0 = 70$~km~s$^{-1}$~Mpc$^{-1}$). Python code using 5-point stencil differentiation is in the Appendix.

\acknowledgments

Several predictions of the broader framework were falsified during development: a cosmic rotation interpretation (ruled out by \citealt{Saadeh2016} at $10^7\times$); a hemispheric density asymmetry (SDSS survey artifact, sign reversal confirmed in all-sky Cosmicflows-3); a Cold Spot polarization signal ($0.003\sigma$ after null test on 1,000 simulated sky patches). These failures informed the present model.

\appendix

\section{Python Implementation}

\begin{verbatim}
import numpy as np
from scipy import integrate

RETURN_FRAC = 0.4

def MD14(z):
    return 0.015*(1+z)**2.7 / (1+((1+z)/2.9)**5.6)

def H_LCDM(z):
    H0_SI = 70*1e3/3.086e22
    return H0_SI*np.sqrt(0.315*(1+z)**3 + 0.685)

def surviving_stellar_mass(z, z_max=10):
    z_arr = np.linspace(z, z_max, 400)
    sfr = MD14(z_arr)
    dt_dz = -1/((1+z_arr)*H_LCDM(z_arr))
    return (1-RETURN_FRAC) * (
        -integrate.simpson(sfr*dt_dz, x=z_arr))

M0 = surviving_stellar_mass(0)
f_prim, gamma = 0.445, 0.272

def net_integral(z, z_max=10):
    z_arr = np.linspace(z, z_max, 300)
    sfr = MD14(z_arr)
    mstar = np.array([surviving_stellar_mass(zi)/M0
                      for zi in z_arr])
    dt = -1/((1+z_arr)*H_LCDM(z_arr))
    net = sfr - gamma*mstar*np.max(sfr)
    return -integrate.simpson(net*dt, x=z_arr)

F0 = net_integral(0)

def rho_DE(z):
    return f_prim + (1-f_prim)*net_integral(z)/F0

def w_eff(z, h=1e-3):
    # 5-point stencil
    rp2 = np.log(rho_DE(z + 2*h))
    rp1 = np.log(rho_DE(z + h))
    rm1 = np.log(rho_DE(z - h))
    rm2 = np.log(rho_DE(z - 2*h))
    dlnrho = (-rp2 + 8*rp1 - 8*rm1 + rm2) / (12*h)
    return -1 + (1+z)/3 * dlnrho
\end{verbatim}

\bibliography{references}
\bibliographystyle{aasjournal}

\end{document}
